\section{Method}
\label{sec:method}
\subsection{Problem and system overview}
Let $\widehat{T}_{g}\in SE(3)$ be a grasp TCP pose estimated from aerial RGB-D
observations in a shared map frame. A ground base pose is
$q=(p_x,p_y,\psi)\in SE(2)$, with $p$ denoting the BUNKER base-link ground-plane
position and $\psi$ its yaw. An aerial viewpoint is $v=(p_v,\psi_v)$ under the
implemented level-hover assumption. The task is to select observations and an
executable ground stance that lead to physical retrieval within a limited number
of accepted sensing windows. We do not solve an optimal joint belief-space
control problem; we use an interpretable one-step sensing surrogate followed by
explicit execution checks.

Figure~\ref{fig:pipeline} shows the loop. Aerial RGB-D initializes the target;
inverse reachability supplies ground-pose support; measured LiDAR endpoints
update environmental evidence; and a task-weighted viewpoint policy requests
further sensing unless a confirmed candidate passes the common execution screen.
The ground robot then navigates, refines the grasp from wrist RGB-D, and performs
the physical operation. All environment fields use aligned $0.10\,\mathrm{m}$
cells in map coordinates. Algorithms consume sensor observations and public
poses/TF, not scene-generator target coordinates. The simulation-specific
localization and contact backends are disclosed in Section~\ref{sec:limitations}.

\begin{figure}[tb]
\centering
\includegraphics[width=\linewidth]{\figdir/figure1_system_pipeline.pdf}
\caption{\textbf{Manipulation-guided sensing and physical handoff.}
Validated exact base footprints project future manipulation requirements into
environment space. Measured MID360 evidence supplies observation deficit and
operational blocking; blue boxes highlight the relevance-weighted sensing path.
The dashed connector denotes workflow order, not an input from deficit to
confirmation: confirmation independently requires measured $H_t$ and $B_t$.
Full-robot screening precedes another view. A ready candidate enters the shared
execution sequence (green), with $\Dexec$ checked after fresh D435 refinement.
Further observation requires remaining budget and a positive admissible score;
failure may terminate the task at any stage. This is a schematic, not a trial
reconstruction.}
\label{fig:pipeline}
\end{figure}

\subsection{Validated manipulation relevance}
\label{sec:relevance}
An RM4D inverse query \cite{rudorfer2025rm4d} produces base candidates
$\{q_i\}$ for $\widehat{T}_g$. The query uses an explicit rigid frame bridge
between the public map and the calibrated robot-reference convention; returned
poses are transformed back to map. The perception target itself is not shifted
for subsequent manipulation. A separate Ground-task map supplies the required
workspace coverage without modifying the reference RM4D algorithm.

For each evaluated candidate, define a validity indicator $g_i$. It is one only
when reachability, inverse kinematics, the validator's collision checks and
minimum joint-limit margin all pass. For the associated joint solution $\theta_i$,
let $m_i=\min_j\{\theta_{ij}-\theta_j^{\min},\theta_j^{\max}-\theta_{ij}\}$.
The continuous quality is
\begin{equation}
 r_i=\begin{cases}
 \operatorname{clip}(m_i/m_{\mathrm{sat}},0,1),&g_i=1,\\
 0,&g_i=0.
 \end{cases}
 \label{eq:quality}
\end{equation}
The gate requires $m_i\ge m_{\min}$ when a valid joint solution exists;
invalid candidates need not have a defined margin. The executed settings use
$m_{\min}=0.01\,\mathrm{rad}$ and
$m_{\mathrm{sat}}=0.5\,\mathrm{rad}$. Numerical IK/FK residuals are acceptance
diagnostics, not manipulation quality. Travel distance and the baseline's final
ranking score do not enter $r_i$.

For a base-grid cell $C$, choose the exact validated winner
\begin{equation}
 i_C^*=\arg\max_{i:p_i\in C,\,g_i=1} r_i,\qquad
 q_C^*=q_{i_C^*},\qquad R(C)=r_{i_C^*}.
 \label{eq:winner}
\end{equation}
Ties retain the first candidate in the fixed evaluated order, independently of
environment evidence. Geometry is anchored at $q_C^*$, not at the grid-cell
center. This avoids introducing an unvalidated displacement merely to rasterize
support. Non-winning candidates are diagnostic and are not reselected as support
because an environment happens to favor them.

The planner validates at most 256 candidates. A cell containing no evaluated
candidate remains \emph{unassessed}; a cell containing evaluations but no valid
one has no validated feasible support in that evaluated set. Neither statement
proves complete continuous-space infeasibility. Thus $R$ is a validated,
partially assessed interest field, not a complete capability map. No smoothing
fills gaps or spreads support to unassessed base poses.

\subsection{Environment evidence and operational footprints}
\label{sec:evidence}
Each accepted LiDAR window supplies endpoints transformed into map using the
sensor-time public TF. For a horizontal ground plane $z_0$, a usable endpoint
within the ground tolerance supplies ground evidence; one at or above the
obstacle-height threshold supplies occupied evidence. The executed tolerances
are $0.02\,\mathrm{m}$ and $0.05\,\mathrm{m}$, respectively. Returns in the
intermediate height band do not create either vote. Evidence is accumulated
per cell per window, not per individual point. Raw occupancy prioritizes an
occupied vote over a free vote in the same window, and historical occupied
evidence is not erased by later ground returns. Two free-vote windows are
required for the raw FREE label when no occupied evidence exists; other cells
are UNKNOWN. Aerial rays are not carved as entire free XY segments: a ray can
pass above an obstacle without observing the ground underneath it.

Let $N_t(x)$ count informative raw evidence windows for cell $x$. The
observation-deficit score is
\begin{equation}
 u_t(x)=\exp[-N_t(x)/\tau],\qquad \tau=2.
 \label{eq:deficit}
\end{equation}
This is a heuristic, not a calibrated probability or entropy. In particular,
FREE does not force $u_t=0$, and a low deficit does not certify collision freedom.

Let $F(q)$ be the padded continuous rectangular base footprint and
\begin{equation}
 K(q)=\{x:\operatorname{cell}(x)\cap F(q)\ne\varnothing\}.
 \label{eq:footprint}
\end{equation}
The padded half-length and half-width are $0.52$ and $0.39\,\mathrm{m}$.
Closed cell/rectangle overlap includes boundary contact, rather than testing
only whether cell centers lie inside the footprint.

Operational evidence distinguishes measured endpoints associated with the
known target, environmental obstacles, and ambiguous returns. Target association
uses runtime perception and geometric consistency, not scene truth. The target
is checked as a perceived, allowance-expanded continuous object against $F(q)$;
a shared coarse cell alone is insufficient to declare target/base collision.
Environmental occupied cells remain conservative blockers. Complete ambiguous
endpoint records are checked using their declared local uncertainty disks
against $F(q)$; missing endpoint history retains the conservative cell blocker.
The executed disk radius is $0.033\,\mathrm{m}$, an engineering allowance
conditional on the public-map TF profile, not a calibrated physical error bound.
Ambiguous returns are not converted to target or free evidence. The target
association allowance, endpoint uncertainty and execution clearance are distinct
declared quantities, not interchangeable safety guarantees.

Denote this operational blocker by $B_t(q)\in\{0,1\}$. Separately, let $H_t(x)$
count windows containing an actual ground endpoint in $x$, even if the same cell
also contains an occupied endpoint elsewhere. Confirmation requires
\begin{equation}
 A_t(q)=\ind[B_t(q)=0]\,
 \ind[F(q)\text{ is inside the grid}]\,
 \prod_{x\in K(q)}\ind[H_t(x)\ge 2].
 \label{eq:confirmation}
\end{equation}
Keeping $H_t$ separate preserves measured ground presence without clearing
occupied evidence. Unknown ground does not block a pose geometrically, but it
does prevent confirmation. $B_t=1$ means blocked by this operational model,
not a proof of global navigation infeasibility. The same assessment is used
for relevance support and exact-candidate handoff.

\subsection{Task relevance in environment space}
\label{sec:projection}
The nominal and operational support projections are
\begin{align}
 M_{\mathrm{nom}}(x)&=\max_{C:x\in K(q_C^*)} R(C),\label{eq:nominal}\\
 M_{\mathrm{op},t}(x)&=\max_{C:x\in K(q_C^*),\,B_t(q_C^*)=0}R(C).
 \label{eq:operational}
\end{align}
Only validated winners participate. Where nominal support exists but all
covering poses are blocked, operational support is zero. Where no nominal
support exists, the implementation retains a distinct no-validated-support
state and excludes that cell from the task sum; it does not reinterpret the
cell as observed or physically unreachable. The task field is
\begin{equation}
 U_{\mathrm{task},t}(x)=u_t(x)M_{\mathrm{op},t}(x).
 \label{eq:task}
\end{equation}
The maximum avoids counting one environment cell repeatedly simply because
several useful base poses cover it. This formulation retains unknown regions
that could support a useful manipulation stance while excluding already blocked
stances. Occupancy does not overwrite $R(C)$. Grid-clipped footprints may
contribute to the in-grid projection but cannot be confirmed by
Eq.~\eqref{eq:confirmation}. Both manipulation coverage and environment
uncertainty remain explicit.

\subsection{Finite-window task-aware viewpoint selection}
\label{sec:nbv}
A cell being inside the LiDAR field of view does not guarantee that a finite
scan window measures it. We therefore use the installed cyclic scan directions
and publisher angular/discard rules to predict ground intersections. For each
scheduled packet-start phase $s\in\mathcal S$, define $h_s(v,x)$ as one if at
least one emitted ray during the next five-second window reaches cell $x$ on
the nominal ground plane, satisfies range limits, and is not occluded. The
acquisition opportunity is
\begin{equation}
 w_t(v,x)=\frac{1}{|\mathcal S|}\sum_{s\in\mathcal S}h_s(v,x).
 \label{eq:opportunity}
\end{equation}
This uniformly averages scheduled starting phases; it is not an empirically
calibrated return probability. Environmental occupied cells and incomplete
ambiguous evidence form assumed-height prisms; complete ambiguous endpoints
form local cylinders; the perceived target uses its oriented box. The
$1\,\mathrm{m}$ prism/cylinder height is an occlusion surrogate, not an estimated
obstacle height. Unknown space does not occlude. Segment tests distinguish
intersection from passing above an occluder. These predictions never increment
evidence counts, and unmodeled interception or pose/scan mismatch can still cause
missed returns.

One additional informative window would reduce Eq.~\eqref{eq:deficit} by
\begin{equation}
 \Delta u_t(x)=(1-e^{-1/\tau})u_t(x)=\alpha u_t(x).
 \label{eq:delta}
\end{equation}
Holding the current manipulation support fixed for one-step prediction gives
\begin{align}
 G_{\mathrm{task}}(v)&=\alpha\sum_xw_t(v,x)U_{\mathrm{task},t}(x),\label{eq:taskgain}\\
 G_{\mathrm{generic}}(v)&=\alpha\sum_xw_t(v,x)u_t(x).
 \label{eq:genericgain}
\end{align}
Equation~\eqref{eq:taskgain} is a predicted observation-gain surrogate: it favors
observable task-relevant deficit mass, not calibrated expected information
gain, guaranteed confirmation, or the probability of task completion.
The current field is recomputed after actual observations.

Both methods use
\begin{equation}
 S_m(v)=G_m(v)-\lambda C(v),\quad
 C(v)=\|p_v-p_{\mathrm{cur}}\|_2+
 \kappa|\operatorname{wrap}(\psi_v-\psi_{\mathrm{cur}})|.
 \label{eq:score}
\end{equation}
The cost is a displacement/yaw proxy, not a measured path length or energy model.
The shared settings are $\lambda=0.25$ and
$\kappa=0.25\,\mathrm{m/rad}$. Candidates include the current pose and a
task-independent local XY lattice with offsets $\{-2,-1,0,1,2\}\,\mathrm{m}$,
fixed current altitude and yaw alternatives, filtered by common bounds. The
tilted LiDAR mount makes yaw relevant despite its panoramic horizontal field
of view. Ties use cost then fixed candidate index. Figure~\ref{fig:comparison}
shows the controlled distinction: only $G_m$'s manipulation weighting differs.

\begin{figure}[tb]
\centering
\includegraphics[width=\linewidth]{\figdir/figure2_method_comparison.pdf}
\caption{\textbf{The controlled difference in observation value.}
\textbf{a}, Generic values observable deficit throughout the environment.
\textbf{b}, Ours additionally weights deficit by operational manipulation support.
Yellow denotes the known target; teal rectangles illustrate useful continuous
base footprints. Coloured outlines distinguish the regions contributing to each
objective schematically, not exact field boundaries. Both policies share finite-scan
opportunity, candidate views, cost, updates, budget and execution.
The spatial cartoons are not measured heatmaps or predictions of a winning view;
$\alpha=1-e^{-1/\tau}$.}
\label{fig:comparison}
\end{figure}

\subsection{Ground execution interface and stopping}
\label{sec:execution}
After each accepted update, exact winners satisfying Eq.~\eqref{eq:confirmation}
are considered in the fixed relevance order. The common execution layer screens
at most four confirmed candidates using the complete base, arm, gripper, target
and post-grasp load models, with bounded observation-pose, IK and approach
searches. A successful screen terminates active perception and supplies the
exact candidate for navigation; the system does not use the field-cell center
or a previously successful stance. If no screen succeeds, the policy requests
its highest-scoring admissible view while budget and positive score permit.
Nonpositive score, no usable view or budget exhaustion can terminate without
handoff. The initial window counts toward the three-window limit.

Screening is only a preview. BUNKER must actually arrive and stop, wrist D435
must provide a fresh usable target estimate, and the arm must establish a
collision-aware refined pregrasp at the measured arrival pose. We call this
executed discovery stage $\Dexec$. Descent, physical closure, lifting and
retention follow with their own checks. The common $12\,\mathrm{mm}$ development
clearance and SIM-specific velocity-feedback mode are execution settings shared
by all applicable comparisons, not method-specific contributions. A reachable
target does not guarantee a usable camera view or a complete collision-free
operation. Failure at any of these stages remains a task failure, not a reason
to redefine success.
